\chapter{Medium GitHub}\label{chap: Medium GitHub}

This manner of using GitHub is quite common, and we won't need an account to be able to proceed.

\begin{itemize}
    \item We wish to download code for our own use.
    \item We may intend to make some changes to suit our needs, but only for ourselves.
    \item We don't intend to upload our changes to the original authors as we are a particular niche case!
\end{itemize}

We can either:

\begin{itemize}
    \item Download the Zip file with the source code, and get a release if necessary, as already discussed in \chapref{chapter: Very Simple GitHub} and \chapref{chapter: Simple GitHub}.
    \item Alternatively, ``clone'' the repository to our PC.
\end{itemize}

\section{Cloning a Repository}\label{sect: Cloning a repository}

Assuming the procedures outlined in \chapref{chapter: Very Simple GitHub} or \chapref{chapter: Simple GitHub} are not enough for us, then proceed as follows to clone the repository and its full change history, onto our PC.

\begin{itemize}
    \item Locate the required repository using any supplied URL or by searching for it.
    \item Click on the ``<> Code'' button as before.
    \item Click ``HTTPS''
    \item Click the icon to copy the repository URL.
    \item Open a terminal session on your PC.
    \item Navigate to wherever you keep your source code, for me it's all kept under \path{SourceCode}.
    \item Within the \path{SourceCode} directory, type \code{git clone } and paste in the URL just coped, to get something like \code{git clone https://github.com/Norman\allowbreak{}Dunbar/QLAssemblyLanguageMagazine.git}. Press return.
    \item The GitHub repository will be cloned to your PC, creating a new sub-directory named \path{QLAssemblyLanguageMagazine} inside your \path{SourceCode} directory.
    \item Now, go play with the code. Make changes, recompile it, have fun!
\end{itemize}

Remember, you have no ability to upload your changes to the original repository.

\section{Updating the Cloned Repository}

You can, from time to time, update your clone of the repository quite simply.

\begin{itemize}
    \item Open a terminal session on your PC.
    \item Navigate to the \path{SourceCode/QLAssemblyLanguageMagazine} sub-directory.
    \item Type the command ``git pull'' and press return.
    \item Any updates made by the original author(s) will be downloaded to your clone, \emph{overwriting what was there before}---it's best not to do this if you have made any important changes!
\end{itemize}